\documentclass[11pt,a4paper]{article}
\usepackage[utf8]{inputenc}
\usepackage[spanish]{babel}
\usepackage{amsmath}
\usepackage{amssymb}
\usepackage{geometry}
\geometry{margin=1in}
\usepackage{listings}
\usepackage{xcolor}
\usepackage{booktabs}
\usepackage{hyperref}
\usepackage{tikz}
\definecolor{codegreen}{rgb}{0,0.6,0}
\definecolor{codegray}{rgb}{0.5,0.5,0.5}
\definecolor{codepurple}{rgb}{0.58,0,0.82}
\definecolor{backcolour}{rgb}{0.95,0.95,0.96}
\lstset{
    language=C++,
    backgroundcolor=\color{backcolour},   
    commentstyle=\color{codegreen},
    keywordstyle=\color{blue}\bfseries,
    numberstyle=\tiny\color{codegray},
    stringstyle=\color{codepurple},
    basicstyle=\ttfamily\small,
    breakatwhitespace=false,         
    breaklines=true,                 
    captionpos=b,                    
    keepspaces=true,                 
    numbers=left,                    
    numbersep=5pt,                  
    showspaces=false,                
    showstringspaces=false,
    showtabs=false,                  
    tabsize=4,
    frame=single,
    rulecolor=\color{gray!30}
}
\begin{document}
\begin{center}
    {\huge \textbf{Principio de Inclusión-Exclusión (PIE)}} \\[0.2cm]
    \vspace{0.5cm}
\end{center}
El \textbf{Principio de Inclusión-Exclusión (PIE)} es una técnica combinatoria fundamental para calcular el tamaño de la unión de múltiples conjuntos finitos. Permite corregir intersecciones reetidas.\par
Sean los conjuntos finitos $C_1, C_2, C_3, \dots, C_N$. El cardinal de su unión se define matemáticamente como:
\begin{equation}
    \left| \bigcup_{i=1}^{N} C_i \right| = \sum_{i=1}^{N} |C_i| - \sum_{1 \leqslant i < j \leqslant N} |C_i \cap C_j| + \sum_{1 \leqslant i < j < k \leqslant N} |C_i \cap C_j \cap C_k| - \dots + (-1)^{N-1} |C_1 \cap C_2 \cap \dots \cap C_N|
\end{equation}
De forma compacta utilizando sumatorias sobre subconjuntos:
\begin{equation}
    \left| \bigcup_{i=1}^{N} C_i \right| = \sum_{\emptyset \neq S \subseteq \{1, \dots, N\}} (-1)^{|S|-1} \left| \bigcap_{i \in S} C_i \right|
\end{equation}
Donde $S$ representa un subconjunto no vacío del índice de nuestros conjuntos, y $|S|$ es la cantidad de elementos en dicho subconjunto.\par
Para $N = 3$ conjuntos, el principio suma los conjuntos individuales, resta las intersecciones dobles (que fueron contadas dos veces) y finalmente vuelve a sumar la intersección triple (que fue restada por completo en el paso anterior).
\begin{center}
\begin{tikzpicture}[scale=1.2, fill opacity=0.35, draw opacity=1]
    \fill[red!80!white] (90:1) circle (1.5);
    \fill[blue!80!white] (210:1) circle (1.5);
    \fill[green!80!white] (330:1) circle (1.5);
    \draw[thick] (90:1) circle (1.5) node [above=1cm, opacity=1] {$C_1$};
    \draw[thick] (210:1) circle (1.5) node [below left=1cm, opacity=1] {$C_2$};
    \draw[thick] (330:1) circle (1.5) node [below right=1cm, opacity=1] {$C_3$};
    \node[opacity=1] at (90:0) {\small $C_1 \cap C_2 \cap C_3$};
\end{tikzpicture}
\end{center}
La definición de ``intersección'' varía drásticamente según la naturaleza del problema matemático:
\begin{itemize}
    \item \textbf{Divisibilidad y Aritmética (Coprimalidad):}
    Si deseamos contar múltiplos comunes de ciertos números, la intersección de ser divisible por $a$ y $b$ equivale a ser divisible por el Mínimo Común Múltiplo:
    \begin{equation}
        \text{Intersección}(a, b) = \text{lcm}(a, b)
    \end{equation}
    \item \textbf{Condiciones Booleanas o Bitmasks:}
    Si las restricciones corresponden a bits que deben estar apagados o encendidos, la intersección se calcula utilizando la operación a nivel de bits AND:
    \begin{equation}
        \text{Intersección}(A, B) = A \ \& \ B
    \end{equation}
    \item \textbf{Caminos en Grillas con Obstáculos:}
    Si $C_i$ es el conjunto de caminos desde el origen al destino que pasan por el obstáculo $i$, la intersección representa caminos que pasan obligatoriamente por un conjunto específico de obstáculos (esto se resuelve típicamente con combinatoria y ordenamiento).
\end{itemize}
\section*{Ejercicios}
\textbf{Enunciado:} Dado un número entero $N$ y tres números primos distintos $p_1, p_2$ y $p_3$, determina cuántos números en el rango $[1, N]$ \textbf{no} son divisibles por ninguno de estos tres primos.\par
Para resolver este problema, definimos los conjuntos de la siguiente manera:
\begin{itemize}
    \item $C_1$: Conjunto de números en $[1, N]$ divisibles por $p_1$. Su tamaño es $\lfloor N / p_1 \rfloor$.
    \item $C_2$: Conjunto de números en $[1, N]$ divisibles por $p_2$. Su tamaño es $\lfloor N / p_2 \rfloor$.
    \item $C_3$: Conjunto de números en $[1, N]$ divisibles por $p_3$. Su tamaño es $\lfloor N / p_3 \rfloor$.
\end{itemize}
La intersección de estos conjuntos representa los números divisibles simultáneamente por sus respectivos primos. Como los números son primos entre sí, la operación de intersección (mínimo común múltiplo) equivale simplemente a la multiplicación de los mismos (por ejemplo, $|C_1 \cap C_2| = \lfloor N / (p_1 \cdot p_2) \rfloor$).\par
\begin{equation}
    \text{Resultado} = |U| - |C_1 \cup C_2 \cup C_3|
\end{equation}
Para calcular la unión $|C_1 \cup C_2 \cup C_3|$ de forma genérica y limpia, utilizamos una máscara de bits (\textit{bitmask}) que evalúa las $2^3 - 1 = 7$ combinaciones posibles de subconjuntos no vacíos.
\begin{lstlisting}
#include <iostream>
#include <vector>
using namespace std;
#define ll long long
int main(){    
    ll n;cin>>n;    
    vector<ll>p(3);
    for(int i=0;i<3;i++)cin>>p[i];
    ll u=0;
    int t=1<<3;
    for (int mask=1;mask<t;mask++){
        ll pr=1;
        int b=0;
        for(int j=0;j<3;j++){
            if((mask>>j)&1) {
                pr*=p[j];
                b++;
            }
        }
        ll c=n/pr;
        if(b%2==1)u+=c;
        else u-=c;
    }
    cout<<n-u<<'\n';
    return 0;
}
\end{lstlisting}
\end{document}